\documentclass[11pt,a4paper]{article}
\usepackage[utf-8]{inputenc}
\usepackage[margin=1in]{geometry}
\usepackage{amsmath}
\usepackage{amssymb}
\usepackage{graphicx}
\usepackage{booktabs}
\usepackage{xcolor}
\usepackage{listings}
\usepackage{fancyhdr}
\usepackage{hyperref}

\pagestyle{fancy}
\fancyhf{}
\rhead{LuxAeterna LSTM Pipeline Report}
\lhead{May 2, 2026}
\cfoot{\thepage}

\title{\textbf{LuxAeterna: Global Webcam Light Quality Prediction}\\
\Large Data Ingestion, Preprocessing, Feature Engineering, Training \& Results}
\author{Automated ML Pipeline}
\date{May 2, 2026}

\begin{document}

\maketitle

\section{Executive Summary}

This report documents the complete machine learning pipeline for predicting ALQS (Air Light Quality Score, 0--100) from global webcam time-series data and weather observations. The pipeline ingests multi-webcam observations, constructs irregular time-series sequences, and trains an LSTM predictor.

\begin{center}
\Large\textbf{\textcolor{red}{RESULT: FAILURE}}\\
Model performance is \textbf{worse than persistence baseline}.\\
No meaningful signal captured beyond autocorrelation.
\end{center}

\section{1. Data Ingestion}

\subsection{1.1 Webcam Discovery}

\begin{itemize}
    \item \textbf{Provider:} Windy v3 API
    \item \textbf{Target:} 1,000 global webcams
    \item \textbf{Method:} Grid-based geographic sweep using 24 seed locations (every 30--60° latitude/longitude)
    \item \textbf{Result:} Successfully cached exactly 1,000 unique webcams
    \item \textbf{API Key:} Required; loaded from \texttt{.env}
\end{itemize}

\subsection{1.2 Data Collection}

\begin{itemize}
    \item \textbf{Duration:} ~13 hours per webcam (approximately 26 timesteps)
    \item \textbf{Interval:} Intended ~30 minutes; actual intervals irregular (15--45 minutes observed)
    \item \textbf{Sampling Rate:} 2 samples per hour target $\times$ 1,000 webcams $\times$ ~13 hours
    \item \textbf{Ingestion Pipeline:} Global background worker; asynchronous batch processing (50 webcams/batch)
    \item \textbf{Image Storage:} Local JPEG files per date/webcam
    \item \textbf{Weather Source:} Open-Meteo API (archive + forecast blend)
    \item \textbf{Output Format:} Multiple Parquet files in \texttt{data/processed/global\_dataset/}
\end{itemize}

\subsection{1.3 Raw Data Statistics}

\begin{table}[h]
\centering
\begin{tabular}{lr}
\toprule
\textbf{Metric} & \textbf{Value} \\
\midrule
Total Rows (Initial Load) & 26,233 \\
Total Parquet Files & ~461 \\
Unique Webcams & 1,000 \\
Rows per Webcam (avg) & 26.2 \\
Temporal Span & 13.0 hours \\
\bottomrule
\end{tabular}
\caption{Raw ingestion dataset overview.}
\end{table}

\subsection{1.4 Data Columns}

\begin{itemize}
    \item \texttt{timestamp}: UTC observation time (datetime64[ns])
    \item \texttt{webcam\_id}: Unique webcam identifier (string)
    \item \texttt{latitude, longitude}: Geographic coordinates (float)
    \item \texttt{image\_path}: Local file path to JPEG (string)
    \item \texttt{alqs}: Target light quality score 0--100 (float)
    \item \texttt{temperature}: Air temperature °C (float)
    \item \texttt{relative\_humidity}: Humidity \% (float)
    \item \texttt{visibility}: Horizontal visibility meters (float, sometimes null)
    \item \texttt{cloud\_cover\_low, cloud\_cover\_mid, cloud\_cover\_high}: Cloud cover \% by layer (float)
    \item \texttt{weather\_code}: WMO weather code (int)
\end{itemize}

\section{2. Preprocessing \& Data Cleaning}

\subsection{2.1 Cleaning Pipeline}

\begin{enumerate}
    \item \textbf{Timestamp Normalization:} Converted all timestamp columns to UTC datetime64. Handled multiple input formats (seconds, milliseconds, datetime strings).
    
    \item \textbf{Missing Value Handling:}
    \begin{itemize}
        \item Dropped rows with missing critical fields: \texttt{alqs, temperature, relative\_humidity, cloud\_cover\_low}
        \item Imputed \texttt{visibility} with dataset median (10,000 meters when no valid median)
        \item Dropped rows with invalid \texttt{webcam\_id, timestamp}
    \end{itemize}
    
    \item \textbf{Coordinate Validation:} Enforced $-90 \leq \text{latitude} \leq 90$ and $-180 \leq \text{longitude} \leq 180$.
    
    \item \textbf{Deduplication:} Removed duplicate (webcam\_id, timestamp) pairs; kept last observation.
    
    \item \textbf{Rows After Cleaning:} 26,233 $\rightarrow$ 26,233 (0 rows dropped---data quality was excellent).
\end{enumerate}

\subsection{2.2 Derived Features}

\begin{itemize}
    \item \textbf{total\_cloud\_cover:} $\text{cloud\_cover\_low} + \text{cloud\_cover\_mid} + \text{cloud\_cover\_high}$
    \item \textbf{delta\_time\_minutes:} Seconds elapsed since previous observation within same webcam, converted to minutes. Median imputation across gaps (30 min assumed).
\end{itemize}

\section{3. Feature Engineering \& Sequence Construction}

\subsection{3.1 Per-Webcam Grouping \& Sorting}

\begin{itemize}
    \item \textbf{Grouping:} Strictly partitioned dataset by \texttt{webcam\_id}. No cross-webcam mixing.
    \item \textbf{Sorting:} Within each group, sorted by \texttt{timestamp} ascending.
    \item \textbf{Irregular Intervals:} Preserved natural time gaps; did NOT resample to fixed frequency.
\end{itemize}

\subsection{3.2 Sliding Window Construction}

\begin{itemize}
    \item \textbf{Window Size:} 12 timesteps (~6 hours at 30-min intervals)
    \item \textbf{Input (X):} Previous 12 observations
    \item \textbf{Target (y):} ALQS at the next timestep (t+1)
    \item \textbf{Features per Timestep:}
    \begin{enumerate}
        \item temperature
        \item relative\_humidity
        \item cloud\_cover\_low
        \item cloud\_cover\_mid
        \item cloud\_cover\_high
        \item weather\_code
        \item visibility
        \item total\_cloud\_cover
        \item delta\_time\_minutes
    \end{enumerate}
    \textbf{Total:} 9 features
\end{itemize}

\subsection{3.3 Sequence Statistics}

\begin{table}[h]
\centering
\begin{tabular}{lr}
\toprule
\textbf{Metric} & \textbf{Count} \\
\midrule
Total Sequences Built & 14,233 \\
Unique Webcams Contributing & 1,000 \\
Sequences Skipped (< 13 timesteps) & 0 \\
Sequences Skipped (missing data) & 0 \\
\bottomrule
\end{tabular}
\caption{Sequence construction summary.}
\end{table}

\subsection{3.4 Feature Scaling}

\begin{itemize}
    \item \textbf{Scaler:} MinMaxScaler (sklearn)
    \item \textbf{Fit:} On all 14,233 sequences combined
    \item \textbf{Range:} [0, 1]
    \item \textbf{Saved:} \texttt{data/processed/feature\_artifacts.joblib}
    \item \textbf{Inverse Transform:} Available for inference
\end{itemize}

\section{4. Train/Validation/Test Split}

\subsection{4.1 Split Strategy}

\textbf{Objective:} Prevent data leakage from same webcam appearing in multiple partitions.

\begin{itemize}
    \item \textbf{Split Basis:} Webcam ID (not random sequences)
    \item \textbf{Ratio:} 70\% train, 15\% validation, 15\% test
    \item \textbf{Allocation:}
    \begin{itemize}
        \item Train webcams: 700 / 1,000 (70.0\%)
        \item Validation webcams: 150 / 1,000 (15.0\%)
        \item Test webcams: 150 / 1,000 (15.0\%)
    \end{itemize}
\end{itemize}

\subsection{4.2 Split Distribution}

\begin{table}[h]
\centering
\begin{tabular}{lrr}
\toprule
\textbf{Partition} & \textbf{Webcams} & \textbf{Sequences} \\
\midrule
Train & 700 & 9,968 \\
Validation & 150 & 2,137 \\
Test & 150 & 2,128 \\
\midrule
\textbf{Total} & \textbf{1,000} & \textbf{14,233} \\
\bottomrule
\end{tabular}
\caption{Train/validation/test split breakdown.}
\end{table}

\section{5. Model Architecture \& Training}

\subsection{5.1 LSTM Architecture}

\begin{lstlisting}[language=Python, basicstyle=\ttfamily\small]
Input: (batch, 12 timesteps, 9 features)
  |
  v
LSTM Layer 1: 128 units, return_sequences=True
  Dropout: 0.2
  |
  v
LSTM Layer 2: 64 units, return_sequences=False
  Dropout: 0.2
  |
  v
Dense Layer: 32 units, ReLU activation
  |
  v
Output Layer: 1 unit, Linear activation (regression)
\end{lstlisting}

\subsection{5.2 Training Configuration}

\begin{table}[h]
\centering
\begin{tabular}{ll}
\toprule
\textbf{Parameter} & \textbf{Value} \\
\midrule
Optimizer & Adam \\
Initial Learning Rate & 0.001 \\
Loss Function & Mean Squared Error (MSE) \\
Batch Size & 64 \\
Max Epochs & 120 \\
Early Stopping Patience & 10 \\
ReduceLROnPlateau Patience & 5 \\
Validation Set Monitor & val\_loss \\
\bottomrule
\end{tabular}
\caption{Training hyperparameters.}
\end{table}

\subsection{5.3 Training Progress}

\begin{itemize}
    \item \textbf{Epochs Completed:} 33 / 120 (early stopped)
    \item \textbf{Learning Rate Schedule:}
    \begin{enumerate}
        \item Epochs 1--20: $\eta = 0.001$
        \item Epochs 21--27: $\eta = 0.0005$ (reduced at epoch 21)
        \item Epochs 28--32: $\eta = 0.00025$ (reduced at epoch 28)
        \item Epoch 33: $\eta = 0.000125$ (early stop triggered)
    \end{enumerate}
    \item \textbf{Best Validation Loss:} 93.40 (epoch 23)
    \item \textbf{Convergence:} Model stabilized after epoch 2; marginal improvement thereafter
\end{itemize}

\section{6. Results \& Evaluation}

\subsection{6.1 Test Set Performance}

\begin{table}[h]
\centering
\begin{tabular}{lrr}
\toprule
\textbf{Metric} & \textbf{LSTM Model} & \textbf{Persistence Baseline} \\
\midrule
Mean Absolute Error (MAE) & 7.88 & 3.21 \\
Mean Squared Error (MSE) & 99.98 & 33.80 \\
RMSE & 9.999 & 5.82 \\
\bottomrule
\end{tabular}
\caption{Test set metrics: LSTM vs. persistence baseline.}
\end{table}

\subsection{6.2 Detailed Evaluation}

\begin{enumerate}
    \item \textbf{Model MAE:} 7.88 ALQS points
    \item \textbf{Baseline MAE:} 3.21 ALQS points (predicting $y_t = y_{t-1}$)
    \item \textbf{MAE Improvement:} $\frac{3.21 - 7.88}{3.21} \times 100 = -145.21\%$ \\
    \quad $\Rightarrow$ \textcolor{red}{\textbf{Model performs 145\% WORSE than baseline}}
    \item \textbf{Model MSE:} 99.98 $(\text{ALQS}^2)$
    \item \textbf{Baseline MSE:} 33.80 $(\text{ALQS}^2)$
    \item \textbf{MSE Ratio:} $\frac{99.98}{33.80} = 2.96\times$ worse
    \item \textbf{Residual Std Dev:} 9.999 ALQS points
\end{enumerate}

\subsection{6.3 Target Achievement}

\begin{table}[h]
\centering
\begin{tabular}{lcc}
\toprule
\textbf{Target} & \textbf{Threshold} & \textbf{Result} \\
\midrule
MAE $< 12$ & 12.0 & $7.88 < 12$ \quad \textcolor{green}{\checkmark PASS} \\
Baseline Improvement $> 20\%$ & 20.0\% & $-145.21\%$ \quad \textcolor{red}{\times FAIL} \\
\bottomrule
\end{tabular}
\caption{Target metrics vs. achievement.}
\end{table}

\section{7. Analysis \& Interpretation}

\subsection{7.1 Why Did the Model Underperform?}

\begin{enumerate}
    \item \textbf{Strong Autocorrelation Signal:}
    The persistence baseline (predicting $y_t = y_{t-1}$) achieves MAE $= 3.21$. This indicates ALQS values are highly correlated within short (30-min) timescales. The LSTM cannot exploit additional weather features to beat this prior.
    
    \item \textbf{Insufficient Temporal Dynamics:}
    With only 26 observations per webcam over 13 hours, the model sees limited state transitions. ALQS likely exhibits slow dynamics (photometric changes across hours/days), not captured by the 12-step (6-hour) window.
    
    \item \textbf{Weak Weather Coupling:}
    Weather features (temperature, cloud cover, visibility) show weak predictive power for ALQS changes. Photometric quality is driven by:
    \begin{itemize}
        \item Solar geometry (time of day, latitude, season) -- \emph{not} in current features
        \item Scene-specific factors (urban, rural, altitude) -- \emph{not} available
        \item Image quality artifacts (lens, ISO, white balance) -- \emph{not} measured
    \end{itemize}
    Current weather proxies are insufficient.
    
    \item \textbf{Limited Training Horizon:}
    Model trained on 9,968 sequences. While substantial, insufficient to learn complex weather--photometry interactions with only 13-hour temporal context per webcam.
    
    \item \textbf{Outlier Sensitivity:}
    High residual std (9.999) suggests the model produces wild predictions on occasional sequences, dragging average error higher.
\end{enumerate}

\subsection{7.2 Data Quality Assessment}

\begin{itemize}
    \item \textbf{Completeness:} Excellent. Zero rows dropped during cleaning. All critical fields present.
    \item \textbf{Consistency:} Good. Geographic coordinates valid. Timestamps normalized cleanly.
    \item \textbf{Representativeness:} Fair. 1,000 webcams globally distributed. However, only 13-hour window per webcam limits temporal dynamics.
    \item \textbf{Signal Strength:} \textcolor{red}{\textbf{WEAK}}. ALQS signal is dominated by intra-hour autocorrelation, not exogenous weather features.
\end{itemize}

\subsection{7.3 Meaningful Signal Captured?}

\textcolor{red}{\textbf{NO.}} The model captures \textbf{no meaningful signal} beyond the trivial persistence baseline. Evidence:

\begin{itemize}
    \item The baseline (simply copying the previous value) is \textbf{2.46$\times$ better in MAE}.
    \item The model's residuals (std $\approx 10$) are large relative to the ALQS range (0--100).
    \item Zero improvement on 20\% MAE reduction target. Target not only missed but inverted.
\end{itemize}

\section{8. Model Quality Assessment}

\subsection{8.1 Verdict: \textcolor{red}{\textbf{USELESS}}}

\begin{center}
\Large
\begin{tabular}{|c|}
\hline
\textbf{Model Status: } \textcolor{red}{\textbf{REJECTED}} \\
\hline
\end{tabular}
\end{center}

\subsection{8.2 Justification}

\begin{enumerate}
    \item \textbf{Worse than Baseline:} The model is 145\% worse than a trivial baseline (predict last value). No practitioner would deploy it.
    
    \item \textbf{No Learned Structure:} The LSTM fails to extract meaningful patterns from weather features. The network is essentially learning to output values close to the input window mean, which does not beat persistence.
    
    \item \textbf{Impractical for Production:} Any system using this model for real-time ALQS prediction should simply use the previous observed value instead.
    
    \item \textbf{Data Regime Too Short:} 13-hour windows are insufficient to capture photometric dynamics. A full 24--72 hour per-webcam dataset is needed to observe meaningful day/night cycles and weather transitions.
\end{enumerate}

\subsection{8.3 Recommendations for Improvement}

\begin{enumerate}
    \item \textbf{Extended Collection Window:} Ingest $\geq$48--72 hours per webcam at 30-min intervals. Observe full day/night cycles and weather transitions.
    
    \item \textbf{Add Solar Features:} Include computed solar elevation, azimuth, and clear-sky index using PyEphem. These are strong ALQS drivers.
    
    \item \textbf{Longer Sequences:} Use 24--48 timesteps (12--24 hours) to capture diurnal patterns.
    
    \item \textbf{Ensemble Baselines:} Compare against multiple baselines (seasonal naive, moving average, exponential smoothing) before claiming LSTM superiority.
    
    \item \textbf{Scene-Specific Normalization:} Stratify by webcam latitude/longitude and image metadata to account for persistent scene differences.
    
    \item \textbf{Explicit Causality:} Lag weather features relative to ALQS (weather at time $t$, ALQS at time $t+\Delta$) to capture true causal delays.
    
    \item \textbf{Attention Mechanism:} Replace standard LSTM with Attention-based architecture to identify which features/timesteps matter most.
\end{enumerate}

\section{9. Conclusion}

The LuxAeterna LSTM pipeline successfully:

\begin{itemize}
    \item[\checkmark] Ingested 1,000 global webcams
    \item[\checkmark] Constructed 14,233 irregular time-series sequences
    \item[\checkmark] Built and trained an LSTM model end-to-end
    \item[\checkmark] Achieved MAE below 12 on test set
\end{itemize}

However, the model:

\begin{itemize}
    \item[\times] Fails to beat trivial persistence baseline (145\% worse)
    \item[\times] Does not capture meaningful photometric--weather signals
    \item[\times] Is unsuitable for production deployment
    \item[\times] Requires substantially more data and feature engineering
\end{itemize}

\textbf{Overall Assessment:} \textcolor{red}{\textbf{FAILURE}}. The current data collection window (13 hours per webcam) is too short to reveal meaningful temporal dynamics. A redesign with 48--72 hour collection per location, solar features, and longer sequences is essential before re-attempting model development.

\newpage

\section{Appendix: Key File Locations}

\begin{itemize}
    \item \textbf{Webcam Cache:} \texttt{data/webcams.json} (1,000 entries)
    \item \textbf{Raw Dataset:} \texttt{data/processed/global\_dataset/*.parquet} (\textasciitilde461 files)
    \item \textbf{Processed Sequences:} \texttt{data/processed/sequence\_dataset.npz}
    \item \textbf{Feature Artifacts:} \texttt{data/processed/feature\_artifacts.joblib}
    \item \textbf{Feature Metadata:} \texttt{data/processed/feature\_metadata.json}
    \item \textbf{Trained Model:} \texttt{models/artifacts/lstm\_predictor.keras}
    \item \textbf{Training Metadata:} \texttt{models/artifacts/lstm\_metadata.json}
\end{itemize}

\end{document}
